\documentclass[11pt,letterpaper]{article}
\usepackage[technical-reference]{cornell-notes}
\usepackage{needspace}

\title{Clause 25: Iterators Library - Cornell Notes}
\author{}
\date{}
\setDocTitle{Clause 25: Iterators Library}
\setDocSubtitle{C++ 2024 Cornell Notes}
\setDocOwner{C++ 2024 Cornell Notes}
\setDocAuthor{}
\setDocDate{}
\setCornellCollection{C++ 2024 Cornell Notes}
\setCornellUnitType{Clause}
\setCornellUnitNumber{25}
\setCornellUnitTitle{Iterators Library}
\setCornellCompanionLabel{C++ Technical Reference}
\hypersetup{pdftitle={Clause 25: Iterators Library - Cornell Notes},pdfsubject={C++ technical study notes},pdfauthor={}}

\begin{document}
\maketitle
\begin{CornellReferenceOrientationBox}
\textbf{Purpose.} Defines iterator concepts, associated types, movement and distance operations, sentinels, adaptors, stream iterators, and range-access customization points.
\par\textbf{Role.} The traversal abstraction connecting containers, ranges, and algorithms.
\par\textbf{Principal dependencies.} Uses the library-wide rules of Clause 16 together with the applicable core-language concepts and supporting library clauses.
\par\textbf{Primary audience.} generic programmers and algorithm authors.
\par\textbf{Major subject groups.} General; Header <iterator> synopsis; Iterator requirements; Iterator primitives; Iterator adaptors; Stream iterators; Range access.
\end{CornellReferenceOrientationBox}
\section{Learning Objectives}
\begin{itemize}[label=\textcolor{CornellReferenceTeal}{$\blacktriangleright$}]
\item Define the governing ideas of General and apply its stated contract at a call site.
\item Identify the governing ideas of Header <iterator> synopsis and apply its stated contract at a call site.
\item Distinguish the governing ideas of Iterator requirements and apply its stated contract at a call site.
\item Explain the governing ideas of Iterator primitives and apply its stated contract at a call site.
\item Trace the governing ideas of Iterator adaptors and apply its stated contract at a call site.
\item Apply the governing ideas of Stream iterators and apply its stated contract at a call site.
\item Analyze the governing ideas of Range access and apply its stated contract at a call site.
\item Evaluate the governing ideas of General and apply its stated contract at a call site.
\item Diagnose the governing ideas of Associated types and apply its stated contract at a call site.
\item Compare the governing ideas of Customization point objects and apply its stated contract at a call site.
\item Verify the governing ideas of Iterator concepts and apply its stated contract at a call site.
\item Relate the governing ideas of C++17 iterator requirements and apply its stated contract at a call site.
\item Classify the governing ideas of Indirect callable requirements and apply its stated contract at a call site.
\item Justify the governing ideas of Common algorithm requirements and apply its stated contract at a call site.
\item Audit the governing ideas of General and apply its stated contract at a call site.
\item Summarize the governing ideas of Standard iterator tags and apply its stated contract at a call site.
\end{itemize}
\section{Normative Structure Map}
\small\begin{longtable}{@{}>{\RaggedRight\arraybackslash}p{0.13\textwidth}>{\RaggedRight\arraybackslash}p{0.27\textwidth}>{\RaggedRight\arraybackslash}p{0.19\textwidth}>{\RaggedRight\arraybackslash}p{0.31\textwidth}@{}}
\toprule\textbf{Subclause} & \textbf{Subject} & \textbf{Rule category} & \textbf{Key dependencies / entities}\\\midrule\endfirsthead
\toprule\textbf{Subclause} & \textbf{Subject} & \textbf{Rule category} & \textbf{Key dependencies / entities}\\\midrule\endhead
\bottomrule\endfoot
25.1 & General & library semantics & Uses the library-wide rules of Clause 16 together with the applicable core-language concepts and supporting library clauses.\\
25.2 & Header <iterator> synopsis & interface / declarations & Uses the library-wide rules of Clause 16 together with the applicable core-language concepts and supporting library clauses.\\
25.3 & Iterator requirements & requirements & General, Associated types, Customization point objects, Iterator concepts, and related subclauses\\
25.4 & Iterator primitives & library semantics & General, Standard iterator tags, Iterator operations, Range iterator operations\\
25.5 & Iterator adaptors & library semantics & Reverse iterators, Insert iterators, Constant iterators and sentinels, Move iterators and sentinels, and related subclauses\\
25.6 & Stream iterators & library semantics & General, Class template istream\_\allowbreak{}iterator, Class template ostream\_\allowbreak{}iterator, Class template istreambuf\_\allowbreak{}iterator, and related subclauses\\
25.7 & Range access & library semantics & Uses the library-wide rules of Clause 16 together with the applicable core-language concepts and supporting library clauses.\\
\end{longtable}\normalsize
\begin{CornellReferenceNormativeBox}A heading maps the clause; it does not replace the detailed conditions. For each operation, read requirements, effects, result, exceptions, complexity, remarks, and notes with their stated normative force.\end{CornellReferenceNormativeBox}
\subsection*{Rule Classification Guide}
\small\begin{longtable}{@{}>{\RaggedRight\arraybackslash}p{0.25\textwidth}>{\RaggedRight\arraybackslash}p{0.67\textwidth}@{}}\toprule\textbf{Label or category} & \textbf{How to read it}\\\midrule\endfirsthead\toprule\textbf{Label or category} & \textbf{How to read it (continued)}\\\midrule\endhead\bottomrule\endfoot
Constraints & Conditions controlling whether a declaration, overload, or specialization participates in the relevant context.\\
Mandates & Requirements checked for the described declaration or specialization; failure has the stated well-formedness consequence.\\
Preconditions & Obligations the caller establishes before evaluation; violating one forfeits the operation's contract.\\
Effects / Returns / Postconditions & Separate the state change, produced result, and state guaranteed after successful completion.\\
Throws / Error conditions & State how failure is reported and which exceptional paths are permitted or required.\\
Complexity & Bounds or exact counts for the operations named by the specification.\\
Remarks / Recommended practice & Remarks refine applicability or participation; recommended practice guides implementations without silently becoming a program requirement.\\
Exposition-only & A device for describing semantics, not a public name on which a program can depend.\\
\end{longtable}\normalsize
\section{Cornell Cue-and-Notes}
\Needspace{8\baselineskip}
\subsection*{25.1 General}
\begin{longtable}{@{}>{\RaggedRight\arraybackslash\columncolor{CornellReferenceCueGray}}p{0.28\textwidth}>{\RaggedRight\arraybackslash}p{0.64\textwidth}@{}}
\rowcolor{CornellReferenceNavy}\color{white}\textbf{Cue / recall prompt} & \color{white}\textbf{Notes}\\\endfirsthead
\rowcolor{CornellReferenceNavy}\color{white}\textbf{Cue / recall prompt} & \color{white}\textbf{Notes (continued)}\\\endhead
\arrayrulecolor{CornellReferenceLineGray}
\textbf{What contract governs General?}\\[-1mm]{\scriptsize\CornellClauseRef{25.1}} & Frames the terminology, applicability, and relationships needed to interpret General; detailed rules remain in the following subclauses.\\\hline
\end{longtable}
\begin{CornellReferenceSummaryBox}General supplies the library semantics portion of this clause. The key review move is to connect its local wording to the clause-wide model, then classify each condition by normative force before relying on it.\end{CornellReferenceSummaryBox}
\Needspace{8\baselineskip}
\subsection*{25.2 Header <iterator> synopsis}
\begin{longtable}{@{}>{\RaggedRight\arraybackslash\columncolor{CornellReferenceCueGray}}p{0.28\textwidth}>{\RaggedRight\arraybackslash}p{0.64\textwidth}@{}}
\rowcolor{CornellReferenceNavy}\color{white}\textbf{Cue / recall prompt} & \color{white}\textbf{Notes}\\\endfirsthead
\rowcolor{CornellReferenceNavy}\color{white}\textbf{Cue / recall prompt} & \color{white}\textbf{Notes (continued)}\\\endhead
\arrayrulecolor{CornellReferenceLineGray}
\textbf{What contract governs Header <iterator> synopsis?}\\[-1mm]{\scriptsize\CornellClauseRef{25.2}} & Catalogs the declarations made available by Header <iterator> synopsis. The synopsis is an interface map; the detailed specifications supply constraints and behavior.\\\hline
\end{longtable}
\begin{CornellReferenceSummaryBox}Header <iterator> synopsis supplies the interface / declarations portion of this clause. The key review move is to connect its local wording to the clause-wide model, then classify each condition by normative force before relying on it.\end{CornellReferenceSummaryBox}
\Needspace{8\baselineskip}
\subsection*{25.3 Iterator requirements}
\begin{longtable}{@{}>{\RaggedRight\arraybackslash\columncolor{CornellReferenceCueGray}}p{0.28\textwidth}>{\RaggedRight\arraybackslash}p{0.64\textwidth}@{}}
\rowcolor{CornellReferenceNavy}\color{white}\textbf{Cue / recall prompt} & \color{white}\textbf{Notes}\\\endfirsthead
\rowcolor{CornellReferenceNavy}\color{white}\textbf{Cue / recall prompt} & \color{white}\textbf{Notes (continued)}\\\endhead
\arrayrulecolor{CornellReferenceLineGray}
\textbf{What contract governs General?}\\[-1mm]{\scriptsize\CornellClauseRef{25.3.1}} & Frames the terminology, applicability, and relationships needed to interpret General; detailed rules remain in the following subclauses.\\\hline
\textbf{What contract governs Associated types?}\\[-1mm]{\scriptsize\CornellClauseRef{25.3.2}} & Specifies the facility family Associated types. Read its declarations together with mandates, preconditions, effects, results, exceptions, complexity, and lifetime rules.\\\hline
\textbf{What contract governs Customization point objects?}\\[-1mm]{\scriptsize\CornellClauseRef{25.3.3}} & Specifies the facility family Customization point objects. Read its declarations together with mandates, preconditions, effects, results, exceptions, complexity, and lifetime rules.\\\hline
\textbf{What contract governs Iterator concepts?}\\[-1mm]{\scriptsize\CornellClauseRef{25.3.4}} & Explains the traversal or view contract for Iterator concepts; prove domain validity, lifetime, sentinel compatibility, and any borrowing or invalidation property.\\\hline
\textbf{Which obligations govern C++17 iterator requirements?}\\[-1mm]{\scriptsize\CornellClauseRef{25.3.5}} & States the admissibility and semantic obligations for C++17 iterator requirements. Separate compile-time participation rules from caller preconditions and runtime effects.\\\hline
\textbf{Which obligations govern Indirect callable requirements?}\\[-1mm]{\scriptsize\CornellClauseRef{25.3.6}} & States the admissibility and semantic obligations for Indirect callable requirements. Separate compile-time participation rules from caller preconditions and runtime effects.\\\hline
\textbf{Which obligations govern Common algorithm requirements?}\\[-1mm]{\scriptsize\CornellClauseRef{25.3.7}} & States the admissibility and semantic obligations for Common algorithm requirements. Separate compile-time participation rules from caller preconditions and runtime effects.\\\hline
\end{longtable}
\begin{CornellReferenceSummaryBox}Iterator requirements supplies the requirements portion of this clause. The key review move is to connect its local wording to the clause-wide model, then classify each condition by normative force before relying on it.\end{CornellReferenceSummaryBox}
\Needspace{5\baselineskip}\paragraph{Detailed coverage ledger.}
\begin{description}[style=nextline,leftmargin=2.8cm,font=\normalfont\bfseries\color{CornellReferenceTeal}]
\item[25.3.2.1] \textbf{Incrementable traits.} Specifies the facility family Incrementable traits. Read its declarations together with mandates, preconditions, effects, results, exceptions, complexity, and lifetime rules.
\item[25.3.2.2] \textbf{Indirectly readable traits.} Specifies the facility family Indirectly readable traits. Read its declarations together with mandates, preconditions, effects, results, exceptions, complexity, and lifetime rules.
\item[25.3.2.3] \textbf{Iterator traits.} Explains the traversal or view contract for Iterator traits; prove domain validity, lifetime, sentinel compatibility, and any borrowing or invalidation property.
\item[25.3.3.1] \textbf{ranges::iter\_\allowbreak{}move.} Explains the traversal or view contract for ranges::iter\_\allowbreak{}move; prove domain validity, lifetime, sentinel compatibility, and any borrowing or invalidation property.
\item[25.3.3.2] \textbf{ranges::iter\_\allowbreak{}swap.} Governs state transfer or exchange for ranges::iter\_\allowbreak{}swap; check self-operation, validity after movement, exception behavior, and invalidation.
\item[25.3.4.1] \textbf{General.} Frames the terminology, applicability, and relationships needed to interpret General; detailed rules remain in the following subclauses.
\item[25.3.4.2] \textbf{Concept indirectly\_\allowbreak{}readable.} Defines or applies the generic requirement Concept indirectly\_\allowbreak{}readable; syntactic satisfaction must be paired with every stated semantic and equality-preservation expectation.
\item[25.3.4.3] \textbf{Concept indirectly\_\allowbreak{}writable.} Defines or applies the generic requirement Concept indirectly\_\allowbreak{}writable; syntactic satisfaction must be paired with every stated semantic and equality-preservation expectation.
\item[25.3.4.4] \textbf{Concept weakly\_\allowbreak{}incrementable.} Defines or applies the generic requirement Concept weakly\_\allowbreak{}incrementable; syntactic satisfaction must be paired with every stated semantic and equality-preservation expectation.
\item[25.3.4.5] \textbf{Concept incrementable.} Defines or applies the generic requirement Concept incrementable; syntactic satisfaction must be paired with every stated semantic and equality-preservation expectation.
\item[25.3.4.6] \textbf{Concept input\_\allowbreak{}or\_\allowbreak{}output\_\allowbreak{}iterator.} Explains the traversal or view contract for Concept input\_\allowbreak{}or\_\allowbreak{}output\_\allowbreak{}iterator; prove domain validity, lifetime, sentinel compatibility, and any borrowing or invalidation property.
\item[25.3.4.7] \textbf{Concept sentinel\_\allowbreak{}for.} Defines or applies the generic requirement Concept sentinel\_\allowbreak{}for; syntactic satisfaction must be paired with every stated semantic and equality-preservation expectation.
\item[25.3.4.8] \textbf{Concept sized\_\allowbreak{}sentinel\_\allowbreak{}for.} Defines or applies the generic requirement Concept sized\_\allowbreak{}sentinel\_\allowbreak{}for; syntactic satisfaction must be paired with every stated semantic and equality-preservation expectation.
\item[25.3.4.9] \textbf{Concept input\_\allowbreak{}iterator.} Explains the traversal or view contract for Concept input\_\allowbreak{}iterator; prove domain validity, lifetime, sentinel compatibility, and any borrowing or invalidation property.
\item[25.3.4.10] \textbf{Concept output\_\allowbreak{}iterator.} Explains the traversal or view contract for Concept output\_\allowbreak{}iterator; prove domain validity, lifetime, sentinel compatibility, and any borrowing or invalidation property.
\item[25.3.4.11] \textbf{Concept forward\_\allowbreak{}iterator.} Explains the traversal or view contract for Concept forward\_\allowbreak{}iterator; prove domain validity, lifetime, sentinel compatibility, and any borrowing or invalidation property.
\item[25.3.4.12] \textbf{Concept bidirectional\_\allowbreak{}iterator.} Explains the traversal or view contract for Concept bidirectional\_\allowbreak{}iterator; prove domain validity, lifetime, sentinel compatibility, and any borrowing or invalidation property.
\item[25.3.4.13] \textbf{Concept random\_\allowbreak{}access\_\allowbreak{}iterator.} Explains the traversal or view contract for Concept random\_\allowbreak{}access\_\allowbreak{}iterator; prove domain validity, lifetime, sentinel compatibility, and any borrowing or invalidation property.
\item[25.3.4.14] \textbf{Concept contiguous\_\allowbreak{}iterator.} Explains the traversal or view contract for Concept contiguous\_\allowbreak{}iterator; prove domain validity, lifetime, sentinel compatibility, and any borrowing or invalidation property.
\item[25.3.5.1] \textbf{General.} Frames the terminology, applicability, and relationships needed to interpret General; detailed rules remain in the following subclauses.
\item[25.3.5.2] \textbf{Cpp17Iterator.} Explains the traversal or view contract for Cpp17Iterator; prove domain validity, lifetime, sentinel compatibility, and any borrowing or invalidation property.
\item[25.3.5.3] \textbf{Input iterators.} Explains the traversal or view contract for Input iterators; prove domain validity, lifetime, sentinel compatibility, and any borrowing or invalidation property.
\item[25.3.5.4] \textbf{Output iterators.} Explains the traversal or view contract for Output iterators; prove domain validity, lifetime, sentinel compatibility, and any borrowing or invalidation property.
\item[25.3.5.5] \textbf{Forward iterators.} Explains the traversal or view contract for Forward iterators; prove domain validity, lifetime, sentinel compatibility, and any borrowing or invalidation property.
\item[25.3.5.6] \textbf{Bidirectional iterators.} Explains the traversal or view contract for Bidirectional iterators; prove domain validity, lifetime, sentinel compatibility, and any borrowing or invalidation property.
\item[25.3.5.7] \textbf{Random access iterators.} Explains the traversal or view contract for Random access iterators; prove domain validity, lifetime, sentinel compatibility, and any borrowing or invalidation property.
\item[25.3.6.1] \textbf{General.} Frames the terminology, applicability, and relationships needed to interpret General; detailed rules remain in the following subclauses.
\item[25.3.6.2] \textbf{Indirect callable traits.} Specifies the facility family Indirect callable traits. Read its declarations together with mandates, preconditions, effects, results, exceptions, complexity, and lifetime rules.
\item[25.3.6.3] \textbf{Indirect callables.} Specifies the facility family Indirect callables. Read its declarations together with mandates, preconditions, effects, results, exceptions, complexity, and lifetime rules.
\item[25.3.6.4] \textbf{Class template projected.} Specifies the facility family Class template projected. Read its declarations together with mandates, preconditions, effects, results, exceptions, complexity, and lifetime rules.
\item[25.3.7.1] \textbf{General.} Frames the terminology, applicability, and relationships needed to interpret General; detailed rules remain in the following subclauses.
\item[25.3.7.2] \textbf{Concept indirectly\_\allowbreak{}movable.} Defines or applies the generic requirement Concept indirectly\_\allowbreak{}movable; syntactic satisfaction must be paired with every stated semantic and equality-preservation expectation.
\item[25.3.7.3] \textbf{Concept indirectly\_\allowbreak{}copyable.} Defines or applies the generic requirement Concept indirectly\_\allowbreak{}copyable; syntactic satisfaction must be paired with every stated semantic and equality-preservation expectation.
\item[25.3.7.4] \textbf{Concept indirectly\_\allowbreak{}swappable.} Governs state transfer or exchange for Concept indirectly\_\allowbreak{}swappable; check self-operation, validity after movement, exception behavior, and invalidation.
\item[25.3.7.5] \textbf{Concept indirectly\_\allowbreak{}comparable.} Defines or applies the generic requirement Concept indirectly\_\allowbreak{}comparable; syntactic satisfaction must be paired with every stated semantic and equality-preservation expectation.
\item[25.3.7.6] \textbf{Concept permutable.} Defines or applies the generic requirement Concept permutable; syntactic satisfaction must be paired with every stated semantic and equality-preservation expectation.
\item[25.3.7.7] \textbf{Concept mergeable.} Defines or applies the generic requirement Concept mergeable; syntactic satisfaction must be paired with every stated semantic and equality-preservation expectation.
\item[25.3.7.8] \textbf{Concept sortable.} Defines or applies the generic requirement Concept sortable; syntactic satisfaction must be paired with every stated semantic and equality-preservation expectation.
\end{description}
\Needspace{8\baselineskip}
\subsection*{25.4 Iterator primitives}
\begin{longtable}{@{}>{\RaggedRight\arraybackslash\columncolor{CornellReferenceCueGray}}p{0.28\textwidth}>{\RaggedRight\arraybackslash}p{0.64\textwidth}@{}}
\rowcolor{CornellReferenceNavy}\color{white}\textbf{Cue / recall prompt} & \color{white}\textbf{Notes}\\\endfirsthead
\rowcolor{CornellReferenceNavy}\color{white}\textbf{Cue / recall prompt} & \color{white}\textbf{Notes (continued)}\\\endhead
\arrayrulecolor{CornellReferenceLineGray}
\textbf{What contract governs General?}\\[-1mm]{\scriptsize\CornellClauseRef{25.4.1}} & Frames the terminology, applicability, and relationships needed to interpret General; detailed rules remain in the following subclauses.\\\hline
\textbf{What contract governs Standard iterator tags?}\\[-1mm]{\scriptsize\CornellClauseRef{25.4.2}} & Explains the traversal or view contract for Standard iterator tags; prove domain validity, lifetime, sentinel compatibility, and any borrowing or invalidation property.\\\hline
\textbf{What contract governs Iterator operations?}\\[-1mm]{\scriptsize\CornellClauseRef{25.4.3}} & Defines the observable result of Iterator operations together with applicable iterator-domain, ordering, complexity, invalidation, and exception conditions.\\\hline
\textbf{What contract governs Range iterator operations?}\\[-1mm]{\scriptsize\CornellClauseRef{25.4.4}} & Defines the observable result of Range iterator operations together with applicable iterator-domain, ordering, complexity, invalidation, and exception conditions.\\\hline
\end{longtable}
\begin{CornellReferenceSummaryBox}Iterator primitives supplies the library semantics portion of this clause. The key review move is to connect its local wording to the clause-wide model, then classify each condition by normative force before relying on it.\end{CornellReferenceSummaryBox}
\Needspace{5\baselineskip}\paragraph{Detailed coverage ledger.}
\begin{description}[style=nextline,leftmargin=2.8cm,font=\normalfont\bfseries\color{CornellReferenceTeal}]
\item[25.4.4.1] \textbf{General.} Frames the terminology, applicability, and relationships needed to interpret General; detailed rules remain in the following subclauses.
\item[25.4.4.2] \textbf{ranges::advance.} Explains the traversal or view contract for ranges::advance; prove domain validity, lifetime, sentinel compatibility, and any borrowing or invalidation property.
\item[25.4.4.3] \textbf{ranges::distance.} Explains the traversal or view contract for ranges::distance; prove domain validity, lifetime, sentinel compatibility, and any borrowing or invalidation property.
\item[25.4.4.4] \textbf{ranges::next.} Explains the traversal or view contract for ranges::next; prove domain validity, lifetime, sentinel compatibility, and any borrowing or invalidation property.
\item[25.4.4.5] \textbf{ranges::prev.} Explains the traversal or view contract for ranges::prev; prove domain validity, lifetime, sentinel compatibility, and any borrowing or invalidation property.
\end{description}
\Needspace{8\baselineskip}
\subsection*{25.5 Iterator adaptors}
\begin{longtable}{@{}>{\RaggedRight\arraybackslash\columncolor{CornellReferenceCueGray}}p{0.28\textwidth}>{\RaggedRight\arraybackslash}p{0.64\textwidth}@{}}
\rowcolor{CornellReferenceNavy}\color{white}\textbf{Cue / recall prompt} & \color{white}\textbf{Notes}\\\endfirsthead
\rowcolor{CornellReferenceNavy}\color{white}\textbf{Cue / recall prompt} & \color{white}\textbf{Notes (continued)}\\\endhead
\arrayrulecolor{CornellReferenceLineGray}
\textbf{What contract governs Reverse iterators?}\\[-1mm]{\scriptsize\CornellClauseRef{25.5.1}} & Explains the traversal or view contract for Reverse iterators; prove domain validity, lifetime, sentinel compatibility, and any borrowing or invalidation property.\\\hline
\textbf{What contract governs Insert iterators?}\\[-1mm]{\scriptsize\CornellClauseRef{25.5.2}} & Explains the traversal or view contract for Insert iterators; prove domain validity, lifetime, sentinel compatibility, and any borrowing or invalidation property.\\\hline
\textbf{What contract governs Constant iterators and sentinels?}\\[-1mm]{\scriptsize\CornellClauseRef{25.5.3}} & Explains the traversal or view contract for Constant iterators and sentinels; prove domain validity, lifetime, sentinel compatibility, and any borrowing or invalidation property.\\\hline
\textbf{What contract governs Move iterators and sentinels?}\\[-1mm]{\scriptsize\CornellClauseRef{25.5.4}} & Explains the traversal or view contract for Move iterators and sentinels; prove domain validity, lifetime, sentinel compatibility, and any borrowing or invalidation property.\\\hline
\textbf{What contract governs Common iterators?}\\[-1mm]{\scriptsize\CornellClauseRef{25.5.5}} & Explains the traversal or view contract for Common iterators; prove domain validity, lifetime, sentinel compatibility, and any borrowing or invalidation property.\\\hline
\textbf{What contract governs Default sentinel?}\\[-1mm]{\scriptsize\CornellClauseRef{25.5.6}} & Specifies the facility family Default sentinel. Read its declarations together with mandates, preconditions, effects, results, exceptions, complexity, and lifetime rules.\\\hline
\textbf{What contract governs Counted iterators?}\\[-1mm]{\scriptsize\CornellClauseRef{25.5.7}} & Explains the traversal or view contract for Counted iterators; prove domain validity, lifetime, sentinel compatibility, and any borrowing or invalidation property.\\\hline
\textbf{What contract governs Unreachable sentinel?}\\[-1mm]{\scriptsize\CornellClauseRef{25.5.8}} & Specifies the facility family Unreachable sentinel. Read its declarations together with mandates, preconditions, effects, results, exceptions, complexity, and lifetime rules.\\\hline
\end{longtable}
\begin{CornellReferenceSummaryBox}Iterator adaptors supplies the library semantics portion of this clause. The key review move is to connect its local wording to the clause-wide model, then classify each condition by normative force before relying on it.\end{CornellReferenceSummaryBox}
\Needspace{5\baselineskip}\paragraph{Detailed coverage ledger.}
\begin{description}[style=nextline,leftmargin=2.8cm,font=\normalfont\bfseries\color{CornellReferenceTeal}]
\item[25.5.1.1] \textbf{General.} Frames the terminology, applicability, and relationships needed to interpret General; detailed rules remain in the following subclauses.
\item[25.5.1.2] \textbf{Class template reverse\_\allowbreak{}iterator.} Explains the traversal or view contract for Class template reverse\_\allowbreak{}iterator; prove domain validity, lifetime, sentinel compatibility, and any borrowing or invalidation property.
\item[25.5.1.3] \textbf{Requirements.} States the admissibility and semantic obligations for Requirements. Separate compile-time participation rules from caller preconditions and runtime effects.
\item[25.5.1.4] \textbf{Construction and assignment.} Governs state transfer or exchange for Construction and assignment; check self-operation, validity after movement, exception behavior, and invalidation.
\item[25.5.1.5] \textbf{Conversion.} Specifies source and destination conditions for Conversion, the resulting type or value category, and any information-loss, pointer, qualification, or lifetime consequence.
\item[25.5.1.6] \textbf{Element access.} Specifies the facility family Element access. Read its declarations together with mandates, preconditions, effects, results, exceptions, complexity, and lifetime rules.
\item[25.5.1.7] \textbf{Navigation.} Specifies the facility family Navigation. Read its declarations together with mandates, preconditions, effects, results, exceptions, complexity, and lifetime rules.
\item[25.5.1.8] \textbf{Comparisons.} Specifies the facility family Comparisons. Read its declarations together with mandates, preconditions, effects, results, exceptions, complexity, and lifetime rules.
\item[25.5.1.9] \textbf{Non-member functions.} Specifies the facility family Non-member functions. Read its declarations together with mandates, preconditions, effects, results, exceptions, complexity, and lifetime rules.
\item[25.5.2.1] \textbf{General.} Frames the terminology, applicability, and relationships needed to interpret General; detailed rules remain in the following subclauses.
\item[25.5.2.2] \textbf{Class template back\_\allowbreak{}insert\_\allowbreak{}iterator.} Explains the traversal or view contract for Class template back\_\allowbreak{}insert\_\allowbreak{}iterator; prove domain validity, lifetime, sentinel compatibility, and any borrowing or invalidation property.
\item[25.5.2.2.1] \textbf{General.} Frames the terminology, applicability, and relationships needed to interpret General; detailed rules remain in the following subclauses.
\item[25.5.2.2.2] \textbf{Operations.} Defines the observable result of Operations together with applicable iterator-domain, ordering, complexity, invalidation, and exception conditions.
\item[25.5.2.2.3] \textbf{back\_\allowbreak{}inserter.} Specifies the facility family back\_\allowbreak{}inserter. Read its declarations together with mandates, preconditions, effects, results, exceptions, complexity, and lifetime rules.
\item[25.5.2.3] \textbf{Class template front\_\allowbreak{}insert\_\allowbreak{}iterator.} Explains the traversal or view contract for Class template front\_\allowbreak{}insert\_\allowbreak{}iterator; prove domain validity, lifetime, sentinel compatibility, and any borrowing or invalidation property.
\item[25.5.2.3.1] \textbf{General.} Frames the terminology, applicability, and relationships needed to interpret General; detailed rules remain in the following subclauses.
\item[25.5.2.3.2] \textbf{Operations.} Defines the observable result of Operations together with applicable iterator-domain, ordering, complexity, invalidation, and exception conditions.
\item[25.5.2.3.3] \textbf{front\_\allowbreak{}inserter.} Specifies the facility family front\_\allowbreak{}inserter. Read its declarations together with mandates, preconditions, effects, results, exceptions, complexity, and lifetime rules.
\item[25.5.2.4] \textbf{Class template insert\_\allowbreak{}iterator.} Explains the traversal or view contract for Class template insert\_\allowbreak{}iterator; prove domain validity, lifetime, sentinel compatibility, and any borrowing or invalidation property.
\item[25.5.2.4.1] \textbf{General.} Frames the terminology, applicability, and relationships needed to interpret General; detailed rules remain in the following subclauses.
\item[25.5.2.4.2] \textbf{Operations.} Defines the observable result of Operations together with applicable iterator-domain, ordering, complexity, invalidation, and exception conditions.
\item[25.5.2.4.3] \textbf{inserter.} Specifies the facility family inserter. Read its declarations together with mandates, preconditions, effects, results, exceptions, complexity, and lifetime rules.
\item[25.5.3.1] \textbf{General.} Frames the terminology, applicability, and relationships needed to interpret General; detailed rules remain in the following subclauses.
\item[25.5.3.2] \textbf{Alias templates.} Specifies the facility family Alias templates. Read its declarations together with mandates, preconditions, effects, results, exceptions, complexity, and lifetime rules.
\item[25.5.3.3] \textbf{Class template basic\_\allowbreak{}const\_\allowbreak{}iterator.} Explains the traversal or view contract for Class template basic\_\allowbreak{}const\_\allowbreak{}iterator; prove domain validity, lifetime, sentinel compatibility, and any borrowing or invalidation property.
\item[25.5.3.4] \textbf{Member types.} Specifies the facility family Member types. Read its declarations together with mandates, preconditions, effects, results, exceptions, complexity, and lifetime rules.
\item[25.5.3.5] \textbf{Operations.} Defines the observable result of Operations together with applicable iterator-domain, ordering, complexity, invalidation, and exception conditions.
\item[25.5.4.1] \textbf{General.} Frames the terminology, applicability, and relationships needed to interpret General; detailed rules remain in the following subclauses.
\item[25.5.4.2] \textbf{Class template move\_\allowbreak{}iterator.} Explains the traversal or view contract for Class template move\_\allowbreak{}iterator; prove domain validity, lifetime, sentinel compatibility, and any borrowing or invalidation property.
\item[25.5.4.3] \textbf{Requirements.} States the admissibility and semantic obligations for Requirements. Separate compile-time participation rules from caller preconditions and runtime effects.
\item[25.5.4.4] \textbf{Construction and assignment.} Governs state transfer or exchange for Construction and assignment; check self-operation, validity after movement, exception behavior, and invalidation.
\item[25.5.4.5] \textbf{Conversion.} Specifies source and destination conditions for Conversion, the resulting type or value category, and any information-loss, pointer, qualification, or lifetime consequence.
\item[25.5.4.6] \textbf{Element access.} Specifies the facility family Element access. Read its declarations together with mandates, preconditions, effects, results, exceptions, complexity, and lifetime rules.
\item[25.5.4.7] \textbf{Navigation.} Specifies the facility family Navigation. Read its declarations together with mandates, preconditions, effects, results, exceptions, complexity, and lifetime rules.
\item[25.5.4.8] \textbf{Comparisons.} Specifies the facility family Comparisons. Read its declarations together with mandates, preconditions, effects, results, exceptions, complexity, and lifetime rules.
\item[25.5.4.9] \textbf{Non-member functions.} Specifies the facility family Non-member functions. Read its declarations together with mandates, preconditions, effects, results, exceptions, complexity, and lifetime rules.
\item[25.5.4.10] \textbf{Class template move\_\allowbreak{}sentinel.} Specifies the facility family Class template move\_\allowbreak{}sentinel. Read its declarations together with mandates, preconditions, effects, results, exceptions, complexity, and lifetime rules.
\item[25.5.4.11] \textbf{Operations.} Defines the observable result of Operations together with applicable iterator-domain, ordering, complexity, invalidation, and exception conditions.
\item[25.5.5.1] \textbf{Class template common\_\allowbreak{}iterator.} Explains the traversal or view contract for Class template common\_\allowbreak{}iterator; prove domain validity, lifetime, sentinel compatibility, and any borrowing or invalidation property.
\item[25.5.5.2] \textbf{Associated types.} Specifies the facility family Associated types. Read its declarations together with mandates, preconditions, effects, results, exceptions, complexity, and lifetime rules.
\item[25.5.5.3] \textbf{Constructors and conversions.} Specifies how Constructors and conversions establishes object state, including participation constraints, initialization effects, and exceptional exits.
\item[25.5.5.4] \textbf{Accessors.} Specifies the facility family Accessors. Read its declarations together with mandates, preconditions, effects, results, exceptions, complexity, and lifetime rules.
\item[25.5.5.5] \textbf{Navigation.} Specifies the facility family Navigation. Read its declarations together with mandates, preconditions, effects, results, exceptions, complexity, and lifetime rules.
\item[25.5.5.6] \textbf{Comparisons.} Specifies the facility family Comparisons. Read its declarations together with mandates, preconditions, effects, results, exceptions, complexity, and lifetime rules.
\item[25.5.5.7] \textbf{Customizations.} Specifies the facility family Customizations. Read its declarations together with mandates, preconditions, effects, results, exceptions, complexity, and lifetime rules.
\item[25.5.7.1] \textbf{Class template counted\_\allowbreak{}iterator.} Explains the traversal or view contract for Class template counted\_\allowbreak{}iterator; prove domain validity, lifetime, sentinel compatibility, and any borrowing or invalidation property.
\item[25.5.7.2] \textbf{Constructors and conversions.} Specifies how Constructors and conversions establishes object state, including participation constraints, initialization effects, and exceptional exits.
\item[25.5.7.3] \textbf{Accessors.} Specifies the facility family Accessors. Read its declarations together with mandates, preconditions, effects, results, exceptions, complexity, and lifetime rules.
\item[25.5.7.4] \textbf{Element access.} Specifies the facility family Element access. Read its declarations together with mandates, preconditions, effects, results, exceptions, complexity, and lifetime rules.
\item[25.5.7.5] \textbf{Navigation.} Specifies the facility family Navigation. Read its declarations together with mandates, preconditions, effects, results, exceptions, complexity, and lifetime rules.
\item[25.5.7.6] \textbf{Comparisons.} Specifies the facility family Comparisons. Read its declarations together with mandates, preconditions, effects, results, exceptions, complexity, and lifetime rules.
\item[25.5.7.7] \textbf{Customizations.} Specifies the facility family Customizations. Read its declarations together with mandates, preconditions, effects, results, exceptions, complexity, and lifetime rules.
\end{description}
\Needspace{8\baselineskip}
\subsection*{25.6 Stream iterators}
\begin{longtable}{@{}>{\RaggedRight\arraybackslash\columncolor{CornellReferenceCueGray}}p{0.28\textwidth}>{\RaggedRight\arraybackslash}p{0.64\textwidth}@{}}
\rowcolor{CornellReferenceNavy}\color{white}\textbf{Cue / recall prompt} & \color{white}\textbf{Notes}\\\endfirsthead
\rowcolor{CornellReferenceNavy}\color{white}\textbf{Cue / recall prompt} & \color{white}\textbf{Notes (continued)}\\\endhead
\arrayrulecolor{CornellReferenceLineGray}
\textbf{What contract governs General?}\\[-1mm]{\scriptsize\CornellClauseRef{25.6.1}} & Frames the terminology, applicability, and relationships needed to interpret General; detailed rules remain in the following subclauses.\\\hline
\textbf{What contract governs Class template istream\_\allowbreak{}iterator?}\\[-1mm]{\scriptsize\CornellClauseRef{25.6.2}} & Explains the traversal or view contract for Class template istream\_\allowbreak{}iterator; prove domain validity, lifetime, sentinel compatibility, and any borrowing or invalidation property.\\\hline
\textbf{What contract governs Class template ostream\_\allowbreak{}iterator?}\\[-1mm]{\scriptsize\CornellClauseRef{25.6.3}} & Explains the traversal or view contract for Class template ostream\_\allowbreak{}iterator; prove domain validity, lifetime, sentinel compatibility, and any borrowing or invalidation property.\\\hline
\textbf{What contract governs Class template istreambuf\_\allowbreak{}iterator?}\\[-1mm]{\scriptsize\CornellClauseRef{25.6.4}} & Explains the traversal or view contract for Class template istreambuf\_\allowbreak{}iterator; prove domain validity, lifetime, sentinel compatibility, and any borrowing or invalidation property.\\\hline
\textbf{What contract governs Class template ostreambuf\_\allowbreak{}iterator?}\\[-1mm]{\scriptsize\CornellClauseRef{25.6.5}} & Explains the traversal or view contract for Class template ostreambuf\_\allowbreak{}iterator; prove domain validity, lifetime, sentinel compatibility, and any borrowing or invalidation property.\\\hline
\end{longtable}
\begin{CornellReferenceSummaryBox}Stream iterators supplies the library semantics portion of this clause. The key review move is to connect its local wording to the clause-wide model, then classify each condition by normative force before relying on it.\end{CornellReferenceSummaryBox}
\Needspace{5\baselineskip}\paragraph{Detailed coverage ledger.}
\begin{description}[style=nextline,leftmargin=2.8cm,font=\normalfont\bfseries\color{CornellReferenceTeal}]
\item[25.6.2.1] \textbf{General.} Frames the terminology, applicability, and relationships needed to interpret General; detailed rules remain in the following subclauses.
\item[25.6.2.2] \textbf{Constructors and destructor.} Specifies how Constructors and destructor establishes object state, including participation constraints, initialization effects, and exceptional exits.
\item[25.6.2.3] \textbf{Operations.} Defines the observable result of Operations together with applicable iterator-domain, ordering, complexity, invalidation, and exception conditions.
\item[25.6.3.1] \textbf{General.} Frames the terminology, applicability, and relationships needed to interpret General; detailed rules remain in the following subclauses.
\item[25.6.3.2] \textbf{Constructors and destructor.} Specifies how Constructors and destructor establishes object state, including participation constraints, initialization effects, and exceptional exits.
\item[25.6.3.3] \textbf{Operations.} Defines the observable result of Operations together with applicable iterator-domain, ordering, complexity, invalidation, and exception conditions.
\item[25.6.4.1] \textbf{General.} Frames the terminology, applicability, and relationships needed to interpret General; detailed rules remain in the following subclauses.
\item[25.6.4.2] \textbf{Class istreambuf\_\allowbreak{}iterator::proxy.} Explains the traversal or view contract for Class istreambuf\_\allowbreak{}iterator::proxy; prove domain validity, lifetime, sentinel compatibility, and any borrowing or invalidation property.
\item[25.6.4.3] \textbf{Constructors.} Specifies how Constructors establishes object state, including participation constraints, initialization effects, and exceptional exits.
\item[25.6.4.4] \textbf{Operations.} Defines the observable result of Operations together with applicable iterator-domain, ordering, complexity, invalidation, and exception conditions.
\item[25.6.5.1] \textbf{General.} Frames the terminology, applicability, and relationships needed to interpret General; detailed rules remain in the following subclauses.
\item[25.6.5.2] \textbf{Constructors.} Specifies how Constructors establishes object state, including participation constraints, initialization effects, and exceptional exits.
\item[25.6.5.3] \textbf{Operations.} Defines the observable result of Operations together with applicable iterator-domain, ordering, complexity, invalidation, and exception conditions.
\end{description}
\Needspace{8\baselineskip}
\subsection*{25.7 Range access}
\begin{longtable}{@{}>{\RaggedRight\arraybackslash\columncolor{CornellReferenceCueGray}}p{0.28\textwidth}>{\RaggedRight\arraybackslash}p{0.64\textwidth}@{}}
\rowcolor{CornellReferenceNavy}\color{white}\textbf{Cue / recall prompt} & \color{white}\textbf{Notes}\\\endfirsthead
\rowcolor{CornellReferenceNavy}\color{white}\textbf{Cue / recall prompt} & \color{white}\textbf{Notes (continued)}\\\endhead
\arrayrulecolor{CornellReferenceLineGray}
\textbf{What contract governs Range access?}\\[-1mm]{\scriptsize\CornellClauseRef{25.7}} & Explains the traversal or view contract for Range access; prove domain validity, lifetime, sentinel compatibility, and any borrowing or invalidation property.\\\hline
\end{longtable}
\begin{CornellReferenceSummaryBox}Range access supplies the library semantics portion of this clause. The key review move is to connect its local wording to the clause-wide model, then classify each condition by normative force before relying on it.\end{CornellReferenceSummaryBox}
\section{Language-Lawyer and Library-Usage Pitfalls}
\begin{CornellReferencePitfallBox}\begin{itemize}[label=\textcolor{CornellReferenceAmber}{$\blacktriangleright$}]
\item Reading a synopsis as the complete behavioral contract.
\item Ignoring mandates, preconditions, exception behavior, or complexity requirements.
\item Using an exposition-only name as though it were public API.
\item Assuming invalidation, ownership, or thread-safety properties that are not stated.
\item Treating successful compilation as proof that semantic requirements and preconditions are met.
\end{itemize}\end{CornellReferencePitfallBox}
\section{Programmer and Implementation Checklist}
\begin{CornellReferenceChecklistBox}\begin{itemize}[label=$\square$]
\item Identify the declaring header or module exposure for each facility.
\item Check template arguments and mandates before evaluating an operation.
\item Establish every precondition at the call site.
\item Record effects, returned value, postconditions, and exceptions separately.
\item Audit iterator, reference, pointer, and object lifetime after mutation.
\item Verify stated complexity and synchronization assumptions.
\item For \CornellClauseRef{25.1}, verify general against the precise rule category and all cited dependencies.
\item For \CornellClauseRef{25.2}, verify header <iterator> synopsis against the precise rule category and all cited dependencies.
\item For \CornellClauseRef{25.3}, verify iterator requirements against the precise rule category and all cited dependencies.
\item For \CornellClauseRef{25.4}, verify iterator primitives against the precise rule category and all cited dependencies.
\end{itemize}\end{CornellReferenceChecklistBox}
\section{Clause Synthesis}
\begin{CornellReferenceOrientationBox}Defines iterator concepts, associated types, movement and distance operations, sentinels, adaptors, stream iterators, and range-access customization points. The traversal abstraction connecting containers, ranges, and algorithms. A sound review distinguishes syntax and participation from runtime preconditions, keeps notes and examples non-mandatory, and follows cross-clause dependencies before making a portability claim.\end{CornellReferenceOrientationBox}
\section{Self-Test Questions}
\begin{enumerate}
\item What role does General play in the clause's overall model?
\item Which normative distinctions must be kept separate when applying Header <iterator> synopsis?
\item How would you audit a program or implementation that depends on Iterator requirements?
\item Compare Iterator primitives with Iterator adaptors; what different question does each answer?
\item What role does Iterator adaptors play in the clause's overall model?
\item Which normative distinctions must be kept separate when applying Stream iterators?
\item How would you audit a program or implementation that depends on Range access?
\item Compare General with Associated types; what different question does each answer?
\item What role does Associated types play in the clause's overall model?
\item Which normative distinctions must be kept separate when applying Customization point objects?
\item How would you audit a program or implementation that depends on Iterator concepts?
\item Compare C++17 iterator requirements with Indirect callable requirements; what different question does each answer?
\item What role does Indirect callable requirements play in the clause's overall model?
\item Which normative distinctions must be kept separate when applying Common algorithm requirements?
\item Scenario: a program relies on an unstated implementation behavior while using General. What must be checked before calling the program portable?
\item Scenario: a program relies on an unstated implementation behavior while using Standard iterator tags. What must be checked before calling the program portable?
\end{enumerate}
\clearpage\section{Answer Key}
\begin{enumerate}
\item General is one part of the clause's library semantics model. Its role is understood by connecting the local rule in 25.6.5.1 to the clause orientation and its dependent concepts.
\item Keep formation and well-formedness separate from runtime preconditions and effects. Also distinguish mandatory wording from notes, implementation choice from undefined behavior, and interface declarations from detailed contracts; see 25.2.
\item Audit Iterator requirements by locating the controlling wording in 25.3, listing inputs and type constraints, proving any preconditions, recording effects and failure behavior, and checking lifetime, invalidation, complexity, and synchronization where applicable.
\item Iterator primitives and Iterator adaptors occupy different steps or facility groups. Analyze each under its own subclause, then follow the explicit dependency between them rather than transferring one set of guarantees to the other; begin at 25.4.
\item Iterator adaptors is one part of the clause's library semantics model. Its role is understood by connecting the local rule in 25.5 to the clause orientation and its dependent concepts.
\item Keep formation and well-formedness separate from runtime preconditions and effects. Also distinguish mandatory wording from notes, implementation choice from undefined behavior, and interface declarations from detailed contracts; see 25.6.
\item Audit Range access by locating the controlling wording in 25.7, listing inputs and type constraints, proving any preconditions, recording effects and failure behavior, and checking lifetime, invalidation, complexity, and synchronization where applicable.
\item General and Associated types occupy different steps or facility groups. Analyze each under its own subclause, then follow the explicit dependency between them rather than transferring one set of guarantees to the other; begin at 25.6.5.1.
\item Associated types is one part of the clause's library semantics model. Its role is understood by connecting the local rule in 25.5.5.2 to the clause orientation and its dependent concepts.
\item Keep formation and well-formedness separate from runtime preconditions and effects. Also distinguish mandatory wording from notes, implementation choice from undefined behavior, and interface declarations from detailed contracts; see 25.3.3.
\item Audit Iterator concepts by locating the controlling wording in 25.3.4, listing inputs and type constraints, proving any preconditions, recording effects and failure behavior, and checking lifetime, invalidation, complexity, and synchronization where applicable.
\item C++17 iterator requirements and Indirect callable requirements occupy different steps or facility groups. Analyze each under its own subclause, then follow the explicit dependency between them rather than transferring one set of guarantees to the other; begin at 25.3.5.
\item Indirect callable requirements is one part of the clause's requirements model. Its role is understood by connecting the local rule in 25.3.6 to the clause orientation and its dependent concepts.
\item Keep formation and well-formedness separate from runtime preconditions and effects. Also distinguish mandatory wording from notes, implementation choice from undefined behavior, and interface declarations from detailed contracts; see 25.3.7.
\item First identify the exact requirement in 25.6.5.1, then classify the relied-on behavior as required, unspecified, implementation-defined, conditionally supported, or undefined. Portability requires satisfying the program obligations and avoiding dependence on a choice not guaranteed in the target environments.
\item First identify the exact requirement in 25.4.2, then classify the relied-on behavior as required, unspecified, implementation-defined, conditionally supported, or undefined. Portability requires satisfying the program obligations and avoiding dependence on a choice not guaranteed in the target environments.
\end{enumerate}
\section{Key-Term Recap}
\begin{longtable}{@{}>{\RaggedRight\arraybackslash}p{0.28\textwidth}>{\RaggedRight\arraybackslash}p{0.64\textwidth}@{}}\toprule\textbf{Term} & \textbf{Working meaning}\\\midrule\endfirsthead\toprule\textbf{Term} & \textbf{Working meaning (continued)}\\\midrule\endhead\bottomrule\endfoot
iterator & A value used to traverse or denote positions in a sequence.\\
sentinel & A value used to mark the end of an iterator's traversal domain.\\
iter\_\allowbreak{}reference\_\allowbreak{}t & The reference-like result of dereferencing an iterator.\\
indirectly\_\allowbreak{}readable & A concept expressing coherent readable value and reference relationships.\\
random\_\allowbreak{}access\_\allowbreak{}iterator & An iterator supporting constant-time indexed movement and ordering.\\
common\_\allowbreak{}iterator & An adaptor unifying distinct iterator and sentinel types.\\
\end{longtable}
\CornellTakeaway{Iterator algorithms are valid only within a proven traversal domain and the exact readable, writable, movement, and sentinel contracts.}
\end{document}
